\vspace{1em}
\noindent\textbf{Lemma: 5.5 (Lemma von Goursat)} \\
Sei $G \subset \mathbb{C}$ offen und $f: G \rightarrow \mathbb{C}$ holomorph. Ist $\Delta \subset G$ ein abgeschlossenes Dreieck, dessen Inneres und Rand ganz in $G$ liegen, so gilt für das Wegintegral über den Rand $\partial\Delta$:
\[
\oint_{\partial\Delta} f(z) \, dz = 0
\]

\vspace{1em}
\noindent\textbf{Beweis:} \\
Wir unterteilen das Ausgangsdreieck $\Delta = \Delta_0$ durch Verbinden der Seitenmitten in 4 kongruente Teildreiecke $\Delta_0^{(1)}, \dots, \Delta_0^{(4)}$. Da sich die Integrale über die inneren Kanten gegenseitig aufheben, gilt:
\[
\oint_{\partial\Delta_0} f(z) \, dz = \sum_{j=1}^{4} \oint_{\partial\Delta_0^{(j)}} f(z) \, dz
\]
Aus der Dreiecksungleichung folgt, dass es mindestens ein Teildreieck geben muss, nennen wir es $\Delta_1$, für das gilt:
\[
\left| \oint_{\partial\Delta_0} f(z) \, dz \right| \le 4 \left| \oint_{\partial\Delta_1} f(z) \, dz \right|
\]
Führt man diese Konstruktion induktiv fort, erhält man eine geschachtelte Folge von Dreiecken $\Delta_0 \supset \Delta_1 \supset \Delta_2 \supset \dots$ mit der Eigenschaft:
\[
\left| \oint_{\partial\Delta_0} f(z) \, dz \right| \le 4^n \left| \oint_{\partial\Delta_n} f(z) \, dz \right|
\]
Nach dem Cantor'schen Durchschnittssatz gibt es genau einen Punkt $z_0 \in \bigcap_{n=0}^{\infty} \Delta_n$. Da $f$ in $z_0$ komplex differenzierbar (also holomorph) ist, lässt sich $f$ linear approximieren:
\[
f(z) = f(z_0) + f'(z_0)(z-z_0) + r(z)(z-z_0)
\]
mit $r(z) \rightarrow 0$ für $z \rightarrow z_0$. Da die Polynome $f(z_0)$ und $f'(z_0)(z-z_0)$ Stammfunktionen besitzen, verschwinden ihre Integrale über $\partial\Delta_n$. Es verbleibt das Integral über den Restterm. Mit der Standard-Längenabschätzung zeigt man, dass dieses für $n \rightarrow \infty$ gegen $0$ konvergiert, woraus die Behauptung folgt. \hfill $\square$

\vspace{1em}
\noindent\textbf{Def: 5.6 (Sterngebiet)} \\
Ein Gebiet $G \subset \mathbb{C}$ heißt \textit{Sterngebiet}, wenn es mindestens einen Punkt $z_0 \in G$ (ein sogenanntes Sternzentrum) gibt, sodass für jeden Punkt $z \in G$ die Verbindungsstrecke $[z_0, z]$ komplett in $G$ liegt.

\vspace{1em}
\noindent\textbf{Beispiel: 5.7} \\
Jedes konvexe Gebiet (z.B. Kreisscheiben, Rechtecke, die Halbebene) ist ein Sterngebiet, da hier die Verbindungsstrecke zwischen \textit{je zwei} beliebigen Punkten komplett im Gebiet liegt. Die geschlitzte Ebene $\mathbb{C} \setminus (-\infty, 0]$ ist ebenfalls ein Sterngebiet (mit Sternzentrum auf der positiven reellen Achse), aber nicht konvex. $\mathbb{C} \setminus \{0\}$ ist hingegen \textit{kein} Sterngebiet.

\vspace{1em}
\noindent\textbf{Satz: 5.8 (Cauchyscher Integralsatz für Sterngebiete)} \\
Sei $G \subset \mathbb{C}$ ein Sterngebiet und $f: G \rightarrow \mathbb{C}$ holomorph. Dann besitzt $f$ auf $G$ eine Stammfunktion. Insbesondere gilt nach Satz 5.4 für jeden geschlossenen stückweisen $C^1$-Weg $\gamma$ in $G$:
\[
\oint_{\gamma} f(z) \, dz = 0
\]

\vspace{1em}
\noindent\textbf{Beweis:} \\
Sei $z_0 \in G$ ein Sternzentrum. Für jedes $z \in G$ definieren wir:
\[
F(z) := \int_{[z_0, z]} f(\zeta) \, d\zeta
\]
Wir wollen zeigen, dass $F'(z) = f(z)$ gilt. Sei $h \in \mathbb{C}$ so klein, dass $z+h \in G$. Da $G$ ein Sterngebiet bezüglich $z_0$ ist, liegt das Dreieck mit den Eckpunkten $z_0, z, z+h$ komplett in $G$. Nach dem Lemma von Goursat (Satz 5.5) ist das Wegintegral über den Rand dieses Dreiecks gleich null:
\[
\int_{[z_0, z]} f(\zeta) \, d\zeta + \int_{[z, z+h]} f(\zeta) \, d\zeta + \int_{[z+h, z_0]} f(\zeta) \, d\zeta = 0
\]
Umgestellt ergibt sich:
\[
F(z+h) - F(z) = \int_{[z, z+h]} f(\zeta) \, d\zeta
\]
Teilt man nun durch $h$ und zieht $f(z)$ ab, kann man den Beweis exakt wie in Satz 5.4 (Teil ii $\implies$ i) zu Ende führen. Die Stetigkeit von $f$ garantiert, dass der Grenzwert für $h \rightarrow 0$ den Wert $f(z)$ annimmt. \hfill $\square$
